% META: {"id": "1SPE-SECDEG-EX-041", "chapitre": "1SPE-SECOND-DEGRE", "type_objet": "exercice", "capacites": ["1SPE-SECDEG-C3", "1SPE-SECDEG-C4", "1SPE-SECDEG-C5"], "capacites_codes": ["C3", "C4", "C5"], "methodes": ["M3", "M4", "M5"], "parcours": 3, "competences": ["chercher", "raisonner", "calculer"], "duree_min": 40, "mode_creation": "ex_nihilo", "sources_inspiration": [], "parametres_sympy": {}, "fichier_tex": "chapitres/1SPE-SECOND-DEGRE/exercices/1SPE-SECDEG-EX-041.tex", "corrige_tex": "chapitres/1SPE-SECOND-DEGRE/corriges/1SPE-SECDEG-CO-041.tex", "status": "generated"}
% BEGIN-VERIFY
% from sympy import *
% X = symbols('X', real=True)  # X = x^2
% x = symbols('x', real=True)
% # Equation bicarree : X^2 - 5X + 4 = 0
% P = X**2 - 5*X + 4
% Delta = 25 - 16
% assert Delta == 9
% X1 = Rational(5 - 3, 2)
% X2 = Rational(5 + 3, 2)
% assert X1 == 1
% assert X2 == 4
% assert P.subs(X, X1) == 0
% assert P.subs(X, X2) == 0
% # x^2 = 1 => x = +/-1
% # x^2 = 4 => x = +/-2
% sols = [-2, -1, 1, 2]
% P_x = x**4 - 5*x**2 + 4
% for s in sols:
%     assert P_x.subs(x, s) == 0
% # Factorisation : (x^2 - 1)(x^2 - 4) = (x-1)(x+1)(x-2)(x+2)
% fact = (x-1)*(x+1)*(x-2)*(x+2)
% assert expand(fact) == P_x
% # Tableau de signes de P_x = (x-2)(x-1)(x+1)(x+2)
% # Racines : -2, -1, 1, 2
% # Signe de P_x sur chaque intervalle
% test_vals = [(-3, 1), (Rational(-3,2), -1), (0, 1), (Rational(3,2), -1), (3, 1)]
% for val, expected_sign in test_vals:
%     val_sign = 1 if P_x.subs(x, val) > 0 else -1
%     assert val_sign == expected_sign
% # Inegalite P_x >= 0 : x in ]-inf,-2] U [-1,1] U [2,+inf[
% assert P_x.subs(x, -3) > 0
% assert P_x.subs(x, Rational(-3,2)) < 0
% assert P_x.subs(x, 0) > 0
% assert P_x.subs(x, Rational(3,2)) < 0
% assert P_x.subs(x, 3) > 0
% # Inegalite P_x <= 0 : x in [-2,-1] U [1,2]
% assert P_x.subs(x, Rational(-3,2)) < 0
% assert P_x.subs(x, Rational(3,2)) < 0
% END-VERIFY
\begin{exercice}{1SPE-SECDEG-EX-041}{3}{40}
On s'intéresse à l'équation \textbf{bicarrée} :
\[
  x^4 - 5x^2 + 4 = 0.
\]

\begin{enumerate}
  \item \textbf{Substitution.} On pose $X = x^2$, avec $X \geq 0$.
  \begin{enumerate}
    \item Réécrire l'équation bicarrée comme une équation du second degré en $X$.
    \item Calculer le discriminant de ce trinôme en $X$ et résoudre l'équation en $X$.
    \item Les valeurs obtenues pour $X$ sont-elles compatibles avec la condition $X \geq 0$ ?
  \end{enumerate}

  \item \textbf{Solutions de l'équation initiale.}
  \begin{enumerate}
    \item Pour chaque valeur de $X$ trouvée, résoudre $x^2 = X$ et lister toutes les solutions réelles de $x^4 - 5x^2 + 4 = 0$.
    \item Vérifier que les solutions trouvées vérifient bien l'équation initiale.
  \end{enumerate}

  \item \textbf{Factorisation.}
  \begin{enumerate}
    \item Écrire $x^4 - 5x^2 + 4$ comme un produit de deux trinômes du second degré en $x$.
    \item Puis factoriser complètement en un produit de quatre facteurs du premier degré.
  \end{enumerate}

  \item \textbf{Tableau de signes.}
  \begin{enumerate}
    \item En utilisant la forme factorisée, dresser le tableau de signes de $x^4 - 5x^2 + 4$ sur $\mathbb{R}$.
    \item Résoudre l'inéquation $x^4 - 5x^2 + 4 \leq 0$.
    \item Résoudre l'inéquation $x^4 - 5x^2 + 4 \geq 0$.
  \end{enumerate}

  \item \textbf{Question ouverte.} On considère maintenant l'équation $x^4 - 5x^2 + k = 0$, où $k$ est un paramètre réel. En posant $X = x^2$, discuter le nombre de solutions réelles de l'équation selon les valeurs de $k$. On distinguera soigneusement les différents cas.
\end{enumerate}
\end{exercice}
